\documentclass[11pt]{article}

\usepackage[margin=1in]{geometry}
\usepackage[T1]{fontenc}
\usepackage[utf8]{inputenc}
\usepackage{lmodern}
\usepackage{amsmath,amssymb}
\usepackage{booktabs}
\usepackage{enumitem}
\usepackage{hyperref}
\usepackage{tabularx}
\usepackage{longtable}
\usepackage{xcolor}

\hypersetup{
  colorlinks=true,
  linkcolor=blue!50!black,
  urlcolor=blue!50!black,
  citecolor=blue!50!black,
  pdftitle={Proof Queue Report for p-vs-np}
}

\title{Proof Queue Report for \texttt{p-vs-np}}
\author{Neuro-symbolic AI Group\\Idiap Research Institute}
\date{March 17, 2026}

\begin{document}
\maketitle

\begin{abstract}
This report captures the current state of the persistent proof queue for \texttt{p-vs-np}.
It records the active conjectures, their theorem pairs, bridge objects, invariants, route lemmas, and the main obligations
currently assigned to each route. The report is intended as an inspectable snapshot of the prove frontier.
\end{abstract}

\section{Queue Snapshot}
\begin{itemize}[leftmargin=*,itemsep=0.25em]
\item Queue version: 2
\item Queue updated (UTC): 2026-03-16T23:02:20.705936+00:00
\item Report generated (UTC): 2026-03-16T23:06:22.367627+00:00
\item Total entries: 6
\item Pending: 6
\item In progress: 0
\item Completed: 0
\end{itemize}

\section{Status Summary}
\begin{center}
\begin{tabular}{lr}
\toprule
Status & Count \\
\midrule
Pending & 6 \\
\bottomrule
\end{tabular}
\end{center}

\section{Pending}
\subsection*{proofq.p-vs-np.001}
\begin{itemize}[leftmargin=*,itemsep=0.25em]
\item Status: pending
\item Priority: 0.779
\item Bridge: \texttt{bridge.001}
\item Route: \texttt{st.def.hardness-magnification-phenomenon-programme-leve} $\rightarrow$ \texttt{st.thm.oliveira-pich-santhanam-magnification-near-state}
\item Candidate object: MCSP witness family
\item Carrier kind: family of instances equipped with a monotone witness measure
\item Invariant: MCSP complexity profile in nonuniform-circuits
\item Conjecture: There exists an admissible MCSP witness family, realized as a family of instances equipped with a monotone witness measure, such that whenever A barely-nontrivial lower bound for an explicit problem Q in a weak model holds for MCSP in nonuniform-circuits, the object organizes the source-side mathematical content while preserving the invariant MCSP complexity profile in nonuniform-circuits and yields major consequences such as NP not in P/poly.
\item Route lemma under test: Prove that a hardness-magnification witness family turns 'A barely-nontrivial lower bound for an explicit problem Q in a weak model' into 'major consequences such as NP not in P/poly' while preserving 'MCSP complexity profile in nonuniform-circuits'.
\item Route synopsis: Move from 'st.def.hardness-magnification-phenomenon-programme-leve' on problem 'MCSP' in model 'nonuniform-circuits' using 'A barely-nontrivial lower bound for an explicit problem Q in a weak model' toward 'st.thm.oliveira-pich-santhanam-magnification-near-state' in model 'nonuniform-circuits' to obtain 'major consequences such as NP not in P/poly'.
\end{itemize}
\paragraph{Theorem Pair}
\textbf{Source.} A hardness magnification theorem has the form:"A barely-nontrivial lower bound for an explicit problem Q in a weak model implies a major lower bound/separation in a strong model."Typically, Q is a meta-problem (e.g. a Gap-MCSP variant) and"major"means consequences such as NP not in P/poly.

\textbf{Target.} In refined regimes, lower bounds for Gap-MCSP slightly beyond known frontiers (e.g. N\textasciicircum{} 1+varepsilon -type circuit lower bounds for suitable promise gaps and input length N=2\textasciicircum{}n ) imply major consequences such as NP not in P/poly.

\paragraph{Proof Obligations}
\begin{itemize}[leftmargin=*,itemsep=0.2em]
\item Prove the source-to-object claim: the source-side content 'A barely-nontrivial lower bound for an explicit problem Q in a weak model' available at 'st.def.hardness-magnification-phenomenon-programme-leve' can be organized into 'MCSP witness family' as a family of instances equipped with a monotone witness measure.
\item Prove the payoff claim: once 'MCSP witness family' carries the route data from 'A barely-nontrivial lower bound for an explicit problem Q in a weak model' toward 'major consequences such as NP not in P/poly', it forces 'NP-not-in-P/poly' at 'st.thm.oliveira-pich-santhanam-magnification-near-state'.
\item Prove the transport claim: every admissible realization of 'MCSP witness family' preserves 'MCSP complexity profile in nonuniform-circuits' while crossing the route nonuniform-circuits.
\item Discharge the route lemma: Prove that a hardness-magnification witness family turns 'A barely-nontrivial lower bound for an explicit problem Q in a weak model' into 'major consequences such as NP not in P/poly' while preserving 'MCSP complexity profile in nonuniform-circuits'.
\end{itemize}
\paragraph{Candidate Lemmas}
\begin{itemize}[leftmargin=*,itemsep=0.2em]
\item Source-to-object lemma: the source-side content 'A barely-nontrivial lower bound for an explicit problem Q in a weak model' at 'st.def.hardness-magnification-phenomenon-programme-leve' yields a realization of 'MCSP witness family' as a family of instances equipped with a monotone witness measure, satisfying admits source-side witness data for 'A barely-nontrivial lower bound for an explicit problem Q in a weak model'; preserves the route invariant 'MCSP complexity profile in nonuniform-circuits'; forces the target-side payoff 'NP-not-in-P/poly'.
\item Transport lemma: every admissible realization of 'MCSP witness family' preserves 'MCSP complexity profile in nonuniform-circuits' while crossing nonuniform-circuits.
\item Payoff lemma: if 'MCSP witness family' is realized, carries 'MCSP complexity profile in nonuniform-circuits', and discharges the route-specific step from 'A barely-nontrivial lower bound for an explicit problem Q in a weak model' to 'major consequences such as NP not in P/poly', then 'NP-not-in-P/poly' holds at 'st.thm.oliveira-pich-santhanam-magnification-near-state'.
\item Route lemma: Prove that a hardness-magnification witness family turns 'A barely-nontrivial lower bound for an explicit problem Q in a weak model' into 'major consequences such as NP not in P/poly' while preserving 'MCSP complexity profile in nonuniform-circuits'.
\end{itemize}
\subsection*{proofq.p-vs-np.002}
\begin{itemize}[leftmargin=*,itemsep=0.25em]
\item Status: pending
\item Priority: 0.758
\item Bridge: \texttt{bridge.002}
\item Route: \texttt{st.def.hardness-magnification-phenomenon-programme-leve} $\rightarrow$ \texttt{st.thm.oliveira-santhanam-hardness-magnification-for-na}
\item Candidate object: MCSP witness family
\item Carrier kind: family of instances equipped with a monotone witness measure
\item Invariant: MCSP complexity profile in nonuniform-circuits
\item Conjecture: There exists an admissible MCSP witness family, realized as a family of instances equipped with a monotone witness measure, such that whenever A barely-nontrivial lower bound for an explicit problem Q in a weak model holds for MCSP in nonuniform-circuits, the object organizes the source-side mathematical content while preserving the invariant MCSP complexity profile in nonuniform-circuits across the transfer nonuniform-circuits $\rightarrow$  formula-size and yields major lower bounds for NP in much stronger models (including superpolynomial formula/circuit lower bounds, and in some parameter... for MCSP in formula-size.
\item Route lemma under test: Prove that a hardness-magnification witness family preserves 'MCSP complexity profile in nonuniform-circuits' while transporting the route from the nonuniform-circuits regime to the formula-size regime strongly enough to derive 'major lower bounds for NP in much stronger models (including superpolynomial formula/circuit lower bounds, and in some parameter...'.
\item Route synopsis: Move from 'st.def.hardness-magnification-phenomenon-programme-leve' on problem 'MCSP' in model 'nonuniform-circuits' using 'A barely-nontrivial lower bound for an explicit problem Q in a weak model' toward 'st.thm.oliveira-santhanam-hardness-magnification-for-na' in model 'formula-size' to obtain 'major lower bounds for NP in much stronger models (including superpolynomial formula/circuit lower bounds, and in some parameter...'.
\end{itemize}
\paragraph{Theorem Pair}
\textbf{Source.} A hardness magnification theorem has the form:"A barely-nontrivial lower bound for an explicit problem Q in a weak model implies a major lower bound/separation in a strong model."Typically, Q is a meta-problem (e.g. a Gap-MCSP variant) and"major"means consequences such as NP not in P/poly.

\textbf{Target.} There exist explicit"meta"problems (including MCSP variants) such that proving certain emph barely-superlinear lower bounds against restricted models (e.g. formulas) would imply major lower bounds for NP in much stronger models (including superpolynomial formula/circuit lower bounds, and in some parameter regimes consequences as strong as NP not in P/poly ).

\paragraph{Proof Obligations}
\begin{itemize}[leftmargin=*,itemsep=0.2em]
\item Prove the source-to-object claim: the source-side content 'A barely-nontrivial lower bound for an explicit problem Q in a weak model' available at 'st.def.hardness-magnification-phenomenon-programme-leve' can be organized into 'MCSP witness family' as a family of instances equipped with a monotone witness measure.
\item Prove the payoff claim: once 'MCSP witness family' carries the route data from 'A barely-nontrivial lower bound for an explicit problem Q in a weak model' toward 'major lower bounds for NP in much stronger models (including superpolynomial formula/circuit lower bounds, and in some parameter...', it forces 'NP-not-in-P/poly' at 'st.thm.oliveira-santhanam-hardness-magnification-for-na'.
\item Prove the transport claim: every admissible realization of 'MCSP witness family' preserves 'MCSP complexity profile in nonuniform-circuits' while crossing the route nonuniform-circuits $\rightarrow$  formula-size.
\item Discharge the route lemma: Prove that a hardness-magnification witness family preserves 'MCSP complexity profile in nonuniform-circuits' while transporting the route from the nonuniform-circuits regime to the formula-size regime strongly enough to derive 'major lower bounds for NP in much stronger models (including superpolynomial formula/circuit lower bounds, and in some parameter...'.
\end{itemize}
\paragraph{Candidate Lemmas}
\begin{itemize}[leftmargin=*,itemsep=0.2em]
\item Source-to-object lemma: the source-side content 'A barely-nontrivial lower bound for an explicit problem Q in a weak model' at 'st.def.hardness-magnification-phenomenon-programme-leve' yields a realization of 'MCSP witness family' as a family of instances equipped with a monotone witness measure, satisfying admits source-side witness data for 'A barely-nontrivial lower bound for an explicit problem Q in a weak model'; preserves the route invariant 'MCSP complexity profile in nonuniform-circuits'; forces the target-side payoff 'NP-not-in-P/poly'.
\item Transport lemma: every admissible realization of 'MCSP witness family' preserves 'MCSP complexity profile in nonuniform-circuits' while crossing nonuniform-circuits $\rightarrow$  formula-size.
\item Payoff lemma: if 'MCSP witness family' is realized, carries 'MCSP complexity profile in nonuniform-circuits', and discharges the route-specific step from 'A barely-nontrivial lower bound for an explicit problem Q in a weak model' to 'major lower bounds for NP in much stronger models (including superpolynomial formula/circuit lower bounds, and in some parameter...', then 'NP-not-in-P/poly' holds at 'st.thm.oliveira-santhanam-hardness-magnification-for-na'.
\item Route lemma: Prove that a hardness-magnification witness family preserves 'MCSP complexity profile in nonuniform-circuits' while transporting the route from the nonuniform-circuits regime to the formula-size regime strongly enough to derive 'major lower bounds for NP in much stronger models (including superpolynomial formula/circuit lower bounds, and in some parameter...'.
\end{itemize}
\subsection*{proofq.p-vs-np.003}
\begin{itemize}[leftmargin=*,itemsep=0.25em]
\item Status: pending
\item Priority: 0.722
\item Bridge: \texttt{bridge.016}
\item Route: \texttt{st.def.hardness-magnification-phenomenon-programme-leve} $\rightarrow$ \texttt{st.thm.atserias-m-uller-simple-general-magnification-fr}
\item Candidate object: MCSP witness family
\item Carrier kind: family of instances equipped with a monotone witness measure
\item Invariant: exponential-lower-bound witness complexity
\item Conjecture: There exists an admissible MCSP witness family, realized as a family of instances equipped with a monotone witness measure, such that whenever A barely-nontrivial lower bound for an explicit problem Q in a weak model holds for MCSP in nonuniform-circuits, the object organizes the source-side mathematical content while preserving the invariant exponential-lower-bound witness complexity and yields There exist general-purpose magnification theorems (not problem-specific) showing that certain mild lower bounds for....
\item Route lemma under test: Prove that a hardness-magnification witness family upgrades the route from major separation consequence to exponential lower bound without losing 'exponential-lower-bound witness complexity'.
\item Route synopsis: Move from 'st.def.hardness-magnification-phenomenon-programme-leve' on problem 'MCSP' in model 'nonuniform-circuits' using 'A barely-nontrivial lower bound for an explicit problem Q in a weak model' toward 'st.thm.atserias-m-uller-simple-general-magnification-fr' to obtain 'There exist general-purpose magnification theorems (not problem-specific) showing that certain mild lower bounds for...'.
\end{itemize}
\paragraph{Theorem Pair}
\textbf{Source.} A hardness magnification theorem has the form:"A barely-nontrivial lower bound for an explicit problem Q in a weak model implies a major lower bound/separation in a strong model."Typically, Q is a meta-problem (e.g. a Gap-MCSP variant) and"major"means consequences such as NP not in P/poly.

\textbf{Target.} There exist general-purpose magnification theorems (not problem-specific) showing that certain mild lower bounds for distinguishers/sparse-code objects can be amplified into stronger circuit/formula lower bounds, achieving magnification thresholds below previously known lower bounds in some settings.

\paragraph{Proof Obligations}
\begin{itemize}[leftmargin=*,itemsep=0.2em]
\item Prove the source-to-object claim: the source-side content 'A barely-nontrivial lower bound for an explicit problem Q in a weak model' available at 'st.def.hardness-magnification-phenomenon-programme-leve' can be organized into 'MCSP witness family' as a family of instances equipped with a monotone witness measure.
\item Prove the payoff claim: once 'MCSP witness family' carries the route data from 'A barely-nontrivial lower bound for an explicit problem Q in a weak model' toward 'There exist general-purpose magnification theorems (not problem-specific) showing that certain mild lower bounds for...', it forces 'exponential-lower-bound' at 'st.thm.atserias-m-uller-simple-general-magnification-fr'.
\item Prove the transport claim: every admissible realization of 'MCSP witness family' preserves 'exponential-lower-bound witness complexity' while crossing the route nonuniform-circuits.
\item Exhibit the model-alignment claim: 'MCSP witness family' gives a faithful translation across nonuniform-circuits.
\end{itemize}
\paragraph{Candidate Lemmas}
\begin{itemize}[leftmargin=*,itemsep=0.2em]
\item Source-to-object lemma: the source-side content 'A barely-nontrivial lower bound for an explicit problem Q in a weak model' at 'st.def.hardness-magnification-phenomenon-programme-leve' yields a realization of 'MCSP witness family' as a family of instances equipped with a monotone witness measure, satisfying admits source-side witness data for 'A barely-nontrivial lower bound for an explicit problem Q in a weak model'; preserves the route invariant 'exponential-lower-bound witness complexity'; aligns the model shift (nonuniform-circuits) with an explicit transport map.
\item Transport lemma: every admissible realization of 'MCSP witness family' preserves 'exponential-lower-bound witness complexity' while crossing nonuniform-circuits.
\item Payoff lemma: if 'MCSP witness family' is realized, carries 'exponential-lower-bound witness complexity', and discharges the route-specific step from 'A barely-nontrivial lower bound for an explicit problem Q in a weak model' to 'There exist general-purpose magnification theorems (not problem-specific) showing that certain mild lower bounds for...', then 'exponential-lower-bound' holds at 'st.thm.atserias-m-uller-simple-general-magnification-fr'.
\item Route lemma: Prove that a hardness-magnification witness family upgrades the route from major separation consequence to exponential lower bound without losing 'exponential-lower-bound witness complexity'.
\end{itemize}
\subsection*{proofq.p-vs-np.004}
\begin{itemize}[leftmargin=*,itemsep=0.25em]
\item Status: pending
\item Priority: 0.715
\item Bridge: \texttt{bridge.007}
\item Route: \texttt{st.thm.karp-many-natural-problems-are-np-complete-karp1} $\rightarrow$ \texttt{st.thm.murray-williams-non-np-hardness-of-mcsp-under-lo}
\item Candidate object: unspecified witness family
\item Carrier kind: family of instances equipped with a monotone witness measure
\item Invariant: consequence schema NP-complete-classification
\item Conjecture: There exists an admissible unspecified witness family, realized as a family of instances equipped with a monotone witness measure, such that whenever NP-complete under polynomial-time many-one reductions holds for the source regime, the object organizes the source-side mathematical content while preserving the invariant consequence schema NP-complete-classification from the source regime to MCSP and yields is not NP-complete. More broadly, the paper develops evidence and conditional consequences suggesting NP-hardness of MCSP under various... for MCSP.
\item Route lemma under test: Prove that a hardness-magnification witness family turns 'NP-complete under polynomial-time many-one reductions' into 'is not NP-complete. More broadly, the paper develops evidence and conditional consequences suggesting NP-hardness of MCSP under various...' while preserving 'consequence schema NP-complete-classification'.
\item Route synopsis: Move from 'st.thm.karp-many-natural-problems-are-np-complete-karp1' using 'NP-complete under polynomial-time many-one reductions' toward 'st.thm.murray-williams-non-np-hardness-of-mcsp-under-lo' on problem 'MCSP' to obtain 'is not NP-complete. More broadly, the paper develops evidence and conditional consequences suggesting NP-hardness of MCSP under various...'.
\end{itemize}
\paragraph{Theorem Pair}
\textbf{Source.} There exist many problems L in combinatorial optimisation/graph theory such that L is NP-complete under polynomial-time many-one reductions.

\textbf{Target.} Under"local"polynomial-time reductions (where each output bit is computable with very small time), MCSP is not NP-complete. More broadly, the paper develops evidence and conditional consequences suggesting NP-hardness of MCSP under various reductions would entail unexpected class separations.

\paragraph{Proof Obligations}
\begin{itemize}[leftmargin=*,itemsep=0.2em]
\item Prove the source-to-object claim: the source-side content 'NP-complete under polynomial-time many-one reductions' available at 'st.thm.karp-many-natural-problems-are-np-complete-karp1' can be organized into 'unspecified witness family' as a family of instances equipped with a monotone witness measure.
\item Prove the payoff claim: once 'unspecified witness family' carries the route data from 'NP-complete under polynomial-time many-one reductions' toward 'is not NP-complete. More broadly, the paper develops evidence and conditional consequences suggesting NP-hardness of MCSP under various...', it forces 'MCSP-hardness-transfer' at 'st.thm.murray-williams-non-np-hardness-of-mcsp-under-lo'.
\item Prove the transport claim: every admissible realization of 'unspecified witness family' preserves 'consequence schema NP-complete-classification' while crossing the route model shift not yet typed.
\item Construct the transfer-law claim: an explicit reduction or transport map sends 'NP-complete under polynomial-time many-one reductions' into 'unspecified witness family' without losing the route to 'MCSP-hardness-transfer'.
\end{itemize}
\paragraph{Candidate Lemmas}
\begin{itemize}[leftmargin=*,itemsep=0.2em]
\item Source-to-object lemma: the source-side content 'NP-complete under polynomial-time many-one reductions' at 'st.thm.karp-many-natural-problems-are-np-complete-karp1' yields a realization of 'unspecified witness family' as a family of instances equipped with a monotone witness measure, satisfying admits source-side witness data for 'NP-complete under polynomial-time many-one reductions'; preserves the route invariant 'consequence schema NP-complete-classification'; provides an explicit transfer law or reduction mechanism on the route.
\item Transport lemma: every admissible realization of 'unspecified witness family' preserves 'consequence schema NP-complete-classification' while crossing model shift not yet typed.
\item Payoff lemma: if 'unspecified witness family' is realized, carries 'consequence schema NP-complete-classification', and discharges the route-specific step from 'NP-complete under polynomial-time many-one reductions' to 'is not NP-complete. More broadly, the paper develops evidence and conditional consequences suggesting NP-hardness of MCSP under various...', then 'MCSP-hardness-transfer' holds at 'st.thm.murray-williams-non-np-hardness-of-mcsp-under-lo'.
\item Reduction lemma: the source-side content 'NP-complete under polynomial-time many-one reductions' reduces to the surrogate object encoded by 'unspecified witness family' without weakening the path to 'MCSP-hardness-transfer'.
\end{itemize}
\subsection*{proofq.p-vs-np.005}
\begin{itemize}[leftmargin=*,itemsep=0.25em]
\item Status: pending
\item Priority: 0.481
\item Bridge: \texttt{bridge.010}
\item Route: \texttt{st.thm.oliveira-pich-santhanam-magnification-near-state} $\rightarrow$ \texttt{st.thm.oliveira-santhanam-hardness-magnification-for-na}
\item Candidate object: MCSP witness family
\item Carrier kind: family of instances equipped with a monotone witness measure
\item Invariant: MCSP complexity profile in nonuniform-circuits
\item Conjecture: There exists an admissible MCSP witness family, realized as a family of instances equipped with a monotone witness measure, such that whenever In refined regimes, lower bounds for Gap-MCSP slightly beyond known frontiers (e.g. N\textasciicircum{} 1+varepsilon -type circuit lower bounds for... holds for MCSP in nonuniform-circuits, the object organizes the source-side mathematical content while preserving the invariant MCSP complexity profile in nonuniform-circuits across the transfer nonuniform-circuits $\rightarrow$  formula-size and yields major lower bounds for NP in much stronger models (including superpolynomial formula/circuit lower bounds, and in some parameter... for MCSP in formula-size.
\item Route lemma under test: Prove that a MCSP witness family preserves 'MCSP complexity profile in nonuniform-circuits' while transporting the route from the nonuniform-circuits regime to the formula-size regime strongly enough to derive 'major lower bounds for NP in much stronger models (including superpolynomial formula/circuit lower bounds, and in some parameter...'.
\item Route synopsis: Move from 'st.thm.oliveira-pich-santhanam-magnification-near-state' on problem 'MCSP' in model 'nonuniform-circuits' using 'In refined regimes, lower bounds for Gap-MCSP slightly beyond known frontiers (e.g. N\textasciicircum{} 1+varepsilon -type circuit lower bounds for...' toward 'st.thm.oliveira-santhanam-hardness-magnification-for-na' in model 'formula-size' to obtain 'major lower bounds for NP in much stronger models (including superpolynomial formula/circuit lower bounds, and in some parameter...'.
\end{itemize}
\paragraph{Theorem Pair}
\textbf{Source.} In refined regimes, lower bounds for Gap-MCSP slightly beyond known frontiers (e.g. N\textasciicircum{} 1+varepsilon -type circuit lower bounds for suitable promise gaps and input length N=2\textasciicircum{}n ) imply major consequences such as NP not in P/poly.

\textbf{Target.} There exist explicit"meta"problems (including MCSP variants) such that proving certain emph barely-superlinear lower bounds against restricted models (e.g. formulas) would imply major lower bounds for NP in much stronger models (including superpolynomial formula/circuit lower bounds, and in some parameter regimes consequences as strong as NP not in P/poly ).
\subsection*{proofq.p-vs-np.006}
\begin{itemize}[leftmargin=*,itemsep=0.25em]
\item Status: pending
\item Priority: 0.472
\item Bridge: \texttt{bridge.003}
\item Route: \texttt{st.def.cook-reckhow-propositional-proof-system} $\rightarrow$ \texttt{st.thm.cook-reckhow-proof-systems-and-cookreckhow1979}
\item Candidate object: TAUT proof progress measure
\item Carrier kind: monotone potential function over proofs or certificates
\item Invariant: TAUT complexity profile in Cook-Reckhow-proof-system
\item Conjecture: There exists an admissible TAUT proof progress measure, realized as a monotone potential function over proofs or certificates, such that whenever A (Cook--Reckhow) propositional proof system is a polynomial-time computable surjection f:bits\textasciicircum{}*twoheadrightarrow TAUT, where TAUT is... holds for TAUT in Cook-Reckhow-proof-system, the object organizes the source-side mathematical content while preserving the invariant TAUT complexity profile in Cook-Reckhow-proof-system and yields There exists a polynomially bounded propositional proof system for TAUT if and only if NP=coNP.
\item Route lemma under test: Prove that a proof certificate family or monotone proof measure turns 'A (Cook--Reckhow) propositional proof system is a polynomial-time computable surjection f:bits\textasciicircum{}*twoheadrightarrow TAUT, where TAUT is...' into 'There exists a polynomially bounded propositional proof system for TAUT if and only if NP=coNP' while preserving 'TAUT complexity profile in Cook-Reckhow-proof-system'.
\item Route synopsis: Move from 'st.def.cook-reckhow-propositional-proof-system' on problem 'TAUT' in model 'Cook-Reckhow-proof-system' using 'A (Cook--Reckhow) propositional proof system is a polynomial-time computable surjection f:bits\textasciicircum{}*twoheadrightarrow TAUT, where TAUT is...' toward 'st.thm.cook-reckhow-proof-systems-and-cookreckhow1979' in model 'Cook-Reckhow-proof-system' to obtain 'There exists a polynomially bounded propositional proof system for TAUT if and only if NP=coNP'.
\end{itemize}
\paragraph{Theorem Pair}
\textbf{Source.} A (Cook--Reckhow) propositional proof system is a polynomial-time computable surjection f:bits\textasciicircum{}*twoheadrightarrow TAUT, where TAUT is the set of propositional tautologies. A string pi is an f -proof of formula varphi if f(pi)=varphi.

\textbf{Target.} There exists a polynomially bounded propositional proof system for TAUT if and only if NP=coNP.

\end{document}
